% ===== tr-07 块4：第 5 章 pp.11--15（原文页 11--15）=====

% [原文 p.11]
事实上，中心纤维是不可约的，而结点奇点很温和（双有理几何术语中的 Kawamata 对数终端，log terminal）。

\subsection{体积渐近与本质骨架}

考虑如上给出的代数 Calabi--Yau 退化族 $X\to S\setminus\{0\}$。我们沿用 [3] 的做法，考虑 $t\to0$ 时 $\int_{X_t}\Omega_t\wedge\overline{\Omega_t}$ 的渐近。在此过程中我们将引入对偶相交复形与本质骨架的概念，它们是编码中心纤维上各除子相交模式的单纯复形。一个重要教训是：Calabi--Yau 流形的测度论极限更接近单纯复形，而不是代数簇；这表明度量极限必定显著不同于由有界次数的射影嵌入所伴随的 Fubini--Study 度量。

% [Tutor Note: 这一段把“度量极限”与“代数极限”对立起来：代数极限是中心纤维 $\mathcal{X}_0$（一个代数簇），而测度论的极限落在单纯复形 $\Delta_{\mathcal{X}}$ 上。原文的 lesson 正是指这种反差，而非断言两者有包含关系。]

一个非常有用的工具是选取一个 snc 模型（参看注 5.2）$\mathcal{X}$ 来填充中心纤维。中心纤维 $\mathcal{X}_0$ 是 snc 除子，其分支为 $E_i$（$i\in I$），我们记 $\mathcal{X}_0=\sum_{i\in I}b_iE_i$。在半稳定 snc 模型这一特殊情形下，$b_i=1$ 对所有 $i\in I$ 成立；这总可以通过有限基变换做到。典范除子 $K_{\mathcal{X}}$ 支撑在 $\mathcal{X}_0$ 上，因为 $K_{\mathcal{X}}$ 有平凡化截面 $\Omega$。我们可以写 $K_{\mathcal{X}}=\sum_i(a_i+b_i-1)E_i$，于是相对对数典范除子为
\[
K^{\log}_{\mathcal{X}/S} := K_{\mathcal{X}} - K_S + \mathcal{X}_{0,\mathrm{red}} - \mathcal{X}_0 = \sum a_iE_i.
\]
把所有的 $a_i$ 平移一个常数 $\kappa$，等价于把 $\Omega$ 乘以 $t^{\kappa}$，这给 $\int_{X_t}\Omega_t\wedge\overline{\Omega_t}$ 带来一个初等因子 $|t|^{2\kappa}$。因此我们总假设 $\min a_i=0$。

% [Tutor Note: 关系式 $K_{\mathcal{X}}=\sum_i(a_i+b_i-1)E_i$ 中的 $b_i$ 来自 $\mathcal{X}_0=\sum b_iE_i$；而 $K^{\log}$ 的定义里出现 $\mathcal{X}_{0,\mathrm{red}}-\mathcal{X}_0=\sum(1-b_i)E_i$，两者相加恰好把 $b_i$ 消掉，才得到 $\sum a_iE_i$。抄写时若漏掉 $\mathcal{X}_{0,\mathrm{red}}-\mathcal{X}_0$ 这一项，会得到 $\sum(a_i+b_i-1)E_i$ 而非 $\sum a_iE_i$。]

引入按 $E_i$ 的相交模式对 $X_t$ 所作的定量分层是有用的。令 $E_J=\cap_{i\in J}E_i$（$J\subset I$），它在非空时不可约。利用 $\mathcal{X}$ 上一个固定的光滑背景 Kähler 度量的距离函数，我们可以写
\[
E^0_J = \{q\in X_t\mid d(q,E_J)\ll1\}\setminus\{q\in X_t\mid d(q,E_{J'})\ll1,\ \text{某个 } J'\supsetneq J\}.
\]

% [Tutor Note: $E^0_J$ 的定义中第二个集合要求 $J'\supsetneq J$（真包含），即“更深的”相交层；这正好把靠近更深层 $E_{J'}$ 的点从 $E_J$ 的管状邻域里剔除。若误写成 $J'\subset J$ 或 $J'\supseteq J$，$E^0_J$ 的含义就变了。]

在 $\emptyset\ne E_J\subset\mathcal{X}$ 附近，我们记 $p=|J|-1$，并在 $\mathcal{X}$ 上引入局部坐标 $z_0,\dots z_n$，使得 $z_0,z_1,\dots,z_p$ 是 $i\in J$ 时 $E_i$ 的定义方程。除子上的条件意味着：在离开更深层的地方，我们可以安排
\[
t=\prod_{0}^{p}z_i^{b_i},
\]
并且
\[
\Omega = u_J\prod_{0}^{p}z_i^{a_i+b_i}d\log z_i\wedge\prod_{p+1}^{n}dz_j
\]
其中 $u_J$ 是某个局部处处非零的全纯函数。按定义，沿 $X_t$ 有 $\Omega=dt\wedge\Omega_t$，于是在 $E^0_J$ 上
\[
\Omega_t = b_0^{-1}u_Jz_0^{a_0}\dots z_p^{a_p}\prod_{1}^{p}d\log z_i\wedge\prod_{p+1}^{n}dz_j.
\]

% [Tutor Note: $\Omega$ 的公式中第一个乘积是 $\prod_{0}^{p}$（从 $z_0$ 起），而 $\Omega_t$ 的公式中相对应是 $\prod_{1}^{p}$（从 $z_1$ 起），且多出因子 $b_0^{-1}$：这是因为 $dt=b_0t\,d\log z_0+\dots$ 把 $d\log z_0$ 吸收进 $dt$。两个乘积的下标起点不同，抄写时极易照抄成同一个范围（600 dpi 裁图核实：p.11，y 585--615 与 y 655--688 PDF 点）。]

% [原文 p.12]
\[
\resizebox{0.97\textwidth}{!}{$\displaystyle \sqrt{-1}^{\,n^2}\Omega_t\wedge\overline{\Omega_t} = |b_0|^{-2}|u_J|^2|z_0|^{2a_0}\dots|z_p|^{2a_p}\prod_{1}^{p}\sqrt{-1}d\log z_i\wedge d\log\bar z_i\wedge\prod_{p+1}^{n}\sqrt{-1}dz_j\wedge d\bar z_j.$}
\]
还应注意，局部方程 $t=\prod_{0}^{p}z_i^{b_i}$ 有 $b_J=\gcd_{i\in J}b_i$ 个解的分支。使用极坐标 $z_i=e^{x_i\log|t|+\sqrt{-1}\theta_i}$（$i\in J$），可以看出 $\int_{E^0_J}\sqrt{-1}^{\,n^2}\Omega_t\wedge\overline{\Omega_t}$ 的量级为 $O(|\log|t||^{l})$，其中 $l=|\{j\in J: a_j=0\}|-1$。

% [Tutor Note: 上式中的指数是 $n^2$ 而不是 $n$：$\sqrt{-1}^{\,n^2}$ 来自 $\Omega_t\wedge\overline{\Omega_t}$ 的归一化因子，在 1600 dpi 裁图下清晰可辨（p.12，x 265--320、y 565--610 PDF 点）。另外 $l=|\{j\in J:a_j=0\}|-1$ 的 $-1$ 不能漏，它来自被 $dt$ 吸收掉的那个方向。]

局部的对数变量 $x_i=\frac{\log|z_i|}{\log|t|}$ 落在单纯形
\[
\Delta_J = \Big\{\sum_{0}^{p}b_ix_i=1,\ 0\le x_i\le1\Big\}
\]
上。这些变量依赖于 $z_i$ 的选取，但由于除子的局部定义方程彼此相差一个处处非零的全纯函数，$x_i$ 的歧义仅为 $O\big(\frac{1}{|\log|t||}\big)$（$0<|t|\ll1$）。采取更全局的观点：这些单纯形如何拼合而成的组合模式，恰好反映除子 $E_i$ 的相交模式。形式上，这一信息被编码在 snc 模型 $\mathcal{X}$ 的\textbf{对偶相交复形} $\Delta_{\mathcal{X}}$ 中：它是这样的多面体复形，其顶点 $v_i$ 对应于 $E_i$，并且当且仅当 $E_J\ne0$ 时，我们给 $i\in J$ 配上以 $v_i$ 为顶点的单纯形 $\Delta_J$。坐标 $x_j$ 随后在 $\Delta_{\mathcal{X}}$ 上定义了一个\emph{分段整仿射结构}。在至多相差上述 $O\big(\frac{1}{|\log|t||}\big)$ 的歧义之下，我们现在有一个\emph{对数映射} $\operatorname{Log}_{\mathcal{X}}:X_t\to\Delta_{\mathcal{X}}$，局部由 $x_i=\frac{\log|z_i|}{\log|t|}$ 描述。于是“混合”空间 $X\sqcup\Delta_{\mathcal{X}}$ 带有一个自然的拓扑，使得点列 $z_k\in X_t$ 收敛到 $x\in\Delta_{\mathcal{X}}$，当且仅当 $t\to0$ 且 $\operatorname{Log}_{\mathcal{X}}(z_k)\to x$。“混合”这一名称指的是代数簇与单纯对象的混合，它比代数族 $X$ 更适于处理测度论极限。

% [Tutor Note: 原文写 “$E_J\ne0$”（而非 $E_J\ne\emptyset$），这是文献中常见的滥用记号，指 $E_J$ 非空；译文照原文。]

% [Tutor Note: 原文在本段把坐标写作 $x_j$（“The coordinates $x_j$ then define...”），与前文用 $x_i$ 描述单个坐标的写法并存；这是原文的记号习惯，译文照原样，未统一。]

该测度还挑出一个特殊的子复形 $Sk(\mathcal{X})$，称为\textbf{本质骨架}，它由 $\Delta_{\mathcal{X}}$ 中顶点对应于满足 $a_i=0$ 的 $E_i$ 的那些单纯形组成。归一化 CY 测度的极限就支撑在这里。$Sk(\mathcal{X})$ 的维数度量了退化 $X$ 的超越程度；它反映在 $\int_{X_t}\Omega_t\wedge\overline{\Omega_t}$ 的增长阶上。该维数可能取到的最大值是 $n$。

在大复结构极限的情形下，$\dim_{\R}Sk(\mathcal{X})=n$。让我们在半稳定 snc 模型中更显式地分析 CY 测度。对于对应于 $Sk(\mathcal{X})$ 中一个 $n$ 维单纯形的 $E_J$，在 $E^0_J$ 上
\begin{equation}
\sqrt{-1}^{\,n^2}\Omega_t\wedge\overline{\Omega_t} = |u_J|^2\prod_{1}^{n}\sqrt{-1}d\log z_i\wedge d\log\bar z_i. \tag{5.7}
\end{equation}
这里 $u_J$ 趋于它在点层 $E_J$ 处的值 $u_J(E_J)$，后者称为 $\Omega$ 的 Poincaré 留数，并且容易看出它与坐标 $z_i$ 的选取无关。由黎曼曲面上的留数定理可以推出：$|u_J(E_J)|^2$ 与这样的 $J$ 的选取无关 [3, Thm. 7.1]。于是，归一化 CY 测度 (5.5) 经推前到 $\Delta_{\mathcal{X}}$ 后，在

% [Tutor Note: (5.7) 的乘积范围是 $\prod_{1}^{n}$（而一般公式中是 $\prod_{1}^{p}$）：此处的 $n$ 维单纯形对应 $p=|J|-1=n$；同时 $a_i=0$ 使 $|z_i|^{2a_i}$ 因子全部消失，故右端只剩 $|u_J|^2$ 与 $d\log z_i\wedge d\log\bar z_i$。]

% [原文 p.13]
$\Delta_J$ 的内部光滑收敛到一个 Lebesgue 测度的常数倍：
\begin{equation}
\resizebox{0.97\textwidth}{!}{$\displaystyle \operatorname{Log}_{\mathcal{X}*}\,d\mu_t = \operatorname{Log}_{\mathcal{X}*}\frac{\Omega_t\wedge\overline{\Omega_t}}{\int_{X_t}\Omega_t\wedge\overline{\Omega_t}}\xrightarrow{t\to0} d\mu_0 := \mathrm{Const}\cdot dx_1\dots dx_n.$} \tag{5.8}
\end{equation}
注意 $dx_1\dots dx_n$ 之所以被典范地定义，是因为 $\Delta_J$ 上存在整仿射结构。作为 $\Delta_{\mathcal{X}}$ 上的测度，极限 $d\mu_0$ 在 $Sk(\mathcal{X})$ 的 $n$ 维面之外的补集上取零测度，因为 $d\mu_t$ 在相应区域上的积分是 $O\big(\frac{1}{|\log|t||}\big)$。式 (5.8) 中的常数与 $J$ 无关，其唯一目的是使 $d\mu_0$ 成为概率测度。

% [Tutor Note: (5.8) 中的 $\operatorname{Log}_{\mathcal{X}*}$ 表示沿对数映射的推前（pushforward），两处都要写；$\xrightarrow{t\to0}$ 是原书排版，不是把 $t$ 代入 $\operatorname{Log}$。$dx_1\dots dx_n$ 不带 $\wedge$ 也与原文一致（视作体积元）。]

\subsection{爆破的影响}

一个主要的告诫是：snc 模型高度不唯一，因为人们总可以沿中心纤维内的某个轨迹爆破一个给定的 snc 模型。关于极化退化的内蕴信息，例如度量与体积测度的极限行为，是不依赖于具体 snc 模型的。

爆破对对偶相交复形与本质骨架的影响已被很好地理解（参看 [49, A.4]）。直观图像如下：
\begin{itemize}
\item 若我们沿一个真包含于 $E_J=\cap_{j\in J}E_j$ 之中、但不包含在任何更小相交层之内的光滑不可约子簇爆破 snc 模型 $\mathcal{X}$，则新的对偶相交复形包含 $\Delta_{\mathcal{X}}$，同时引入一个新顶点，并在 $\Delta_{\mathcal{X}}$ 的某些面之上引入一些新的“翼”。本质骨架保持不变。
\item 若我们沿某个 $E_J=\cap_{j\in J}E_j$ 爆破 snc 模型 $\mathcal{X}$，则新的对偶相交复形是 $\Delta_{\mathcal{X}}$ 的一个细分。若对应于 $E_J$ 的单纯形 $\Delta_J$ 是本质骨架 $Sk(\mathcal{X})$ 的一个面，则 $Sk(\mathcal{X})$ 上诱导出一个细分；否则本质骨架保持不变。
\end{itemize}

虽然 $\Delta_{\mathcal{X}}$ 在爆破下一般会变大，但本质骨架 $Sk(\mathcal{X})$ 只是被细分，而它的分段仿射结构（特别是它的同胚型）是双有理不变量，记作 $Sk(X)$。双有理不变性并不令人意外：本质骨架是 $X_t$ 的测度论极限，而这一性质并不依赖于模型的选取。

% [Tutor Note: 原文这里同时使用了两种记号：$Sk(\mathcal{X})$ 指某个 snc 模型上的骨架复形，而 $Sk(X)$ 指其（只被细分、同胚型不变的）双有理不变量。二者不能混用：$\Delta_{\mathcal{X}}$ 依赖模型，$Sk(X)$ 不依赖。]

\subsection{Kontsevich--Soibelman 猜想}

为了从 SYZ 图景中提取极限信息，Kontsevich 与 Soibelman [48][49] 提出了如下图景。设 $X_t$ 是大复结构极限附近 Calabi--Yau 流形的极化退化，则

% [原文 p.14]
\begin{itemize}
\item （Kähler 类归一化）重新标度后的度量 $\omega_{CY,t}\in\frac{1}{|\log|t||}c_1(\mathcal{O}(1))$ 具有非平凡的有限直径 Gromov--Hausdorff 子列极限。
\item 在本质骨架 $Sk(X)$ 上，除开一个 Hausdorff 余维二的奇异子集之外存在一个\emph{仿射结构}。在光滑轨迹上，存在一个黎曼度量，它是实 Monge--Ampère 方程局部解的 Hessian。这个度量与 Gromov--Hausdorff 极限一致，而该极限按猜想与子列的选取无关。
\item （拓扑）在严格 Calabi--Yau 条件 $h^{p,0}(X_t)=0$（$p<n$）之下，本质骨架同胚于 $S^n$。
\end{itemize}

其启发式想法是：本质骨架应当是假想的 SYZ 纤维化的底，而 SYZ 纤维化至少在一般区域中由对数映射近似。我们简要评论 Kontsevich--Soibelman 猜想的现状（另参看 [53, 第 4.5 节]）：
\begin{itemize}
\item 与小 $t$ 无关的一致直径估计
\[
C^{-1}\le\mathrm{diam}(X,\omega_{CY,t})\le C
\]
最近在与 Tosatti 的合作工作 [56] 中建立，主要使用黎曼几何方法。它与 Gromov 紧性一起证明了 Kähler 类归一化这一预测。
\item 关于底上除开余维二之外存在度量上更优的仿射结构这一预测，属于 SYZ 猜想的一般常识（参看第 5.2.4 节），尽管证据不足。一般而言，$Sk(X)$ 上的单纯复形结构诱导出一个分段仿射结构，而要把它改进为除开余维二之外的仿射结构，需要高度不平凡的选择。作者不知道下面这个初等问题是否有完全令人满意的答案：
\begin{question}{5.5}
给定大复结构极限附近的一族无特殊对称性的四次 K3 曲面，我们如何确定 $Sk(X)$ 上奇异点的位置？如何显式写出正则轨迹上的仿射结构？
\end{question}
\item 本质骨架的拓扑是双有理几何中一个活跃的研究课题。$Sk(X)$ 与 $S^n$ 之间的同胚在许多例子中已被验证。一般地，已知 [62][63] $Sk(X)$ 是一个“伪流形”，其有理同调群与 $S^n$ 一致，其基本群的 profinite 完备化是平凡的，但实际的同胚型仍难以捉摸，据说需要诉诸 Poincaré 猜想。
\end{itemize}

% [Tutor Note: 这里的 $h^{p,0}(X_t)=0$（$p<n$）即“严格 Calabi--Yau”条件：它排除全纯 $p$-形式，使 $X_t$ 的 Hodge 结构尽可能“刚性”；这是拓扑预测 $Sk(X)\cong S^n$ 的假设之一。注意下标是 $p<n$ 而非 $p\le n$。]

% [Tutor Note: 一致直径估计 $C^{-1}\le\operatorname{diam}(X,\omega_{CY,t})\le C$ 中的常数 $C$ 不依赖 $t$；这正是“非坍缩”在度量层面的定量体现，也是 Gromov 紧性可用的前提。]

% [原文 p.15]
\begin{remark}{5.3}
Gromov--Hausdorff 收敛是赋予弱极限以意义的标准化方式，但也可能存在其他弱概念。Kontsevich 与 Soibelman [49][48] 旨在把非阿基米德几何确立为研究 Calabi--Yau 流形之极限及其对镜像对称之推论的一个合适框架。Kontsevich 与 Tscinkel [50] 开创了模仿 Kähler 几何来建立非阿基米德势理论的尝试，这一任务由 Boucksom 等人 [4][3][6][5] 大大推进（参看第 5.5 节）。Kontsevich 与 Soibelman 也许早在任何严格定义出现之前就已预见到：Calabi--Yau 度量应当在某种势论意义上收敛到一个非阿基米德对象，而极限信息应当可以纯粹由本质骨架上的数据读出。
\end{remark}

% [Tutor Note: 原文 p.15 把 Tschinkel 拼作 “Tscinkel”（漏掉字母 h，900 dpi 裁图核实：p.15，y 170--190 PDF 点）；译文按原文拼写保留为 Tscinkel，读者宜对应到文献 [50] 的 Kontsevich--Tschinkel 工作。]

\section{分析基础}

\subsection{Yau 对 Calabi 猜想的解决}

Kähler 几何的一块基石是 Yau 对 Calabi 猜想的著名证明，它蕴含

\begin{theorem}{5.4.1}
[77] 带有处处非零全纯体积形式 $\Omega$ 的紧 Kähler 流形 $X$，在任意给定的 Kähler 类中都容许唯一的 Calabi--Yau 度量 $(g,\omega,J,\Omega)$。
\end{theorem}

我们现在对 Yau 证明（的现代观点）作一个非常粗略的叙述。固定给定类中的一个背景 Kähler 度量 $\omega$，任务是找到一个 Kähler 势 $\phi$，使得对任意给定的、满足上同调约束
\[
\int_X e^f\omega^n=\int_X\omega^n
\]
的密度函数 $e^f$，$\phi$ 解\textbf{复 Monge--Ampère 方程}
\begin{equation}
\omega_\phi^n=(\omega+dd^c\phi)^n=e^f\omega^n,\qquad \int_X\phi\omega^n=0, \tag{5.9}
\end{equation}
特别地，取 $e^f\omega^n=\mathrm{const}\cdot\Omega\wedge\overline{\Omega}$ 便得到 Calabi--Yau 度量。唯一性由一个简单的分部积分论证得到。

% [Tutor Note: (5.9) 中的归一化条件 $\int_X\phi\omega^n=0$ 用来固定 $\phi$ 的加性常数；没有它，$\omega+dd^c\phi$ 不变但 $\phi$ 不唯一。式 (5.9) 同时写出方程与归一化，二者不可拆开抄。]

存在性证明采用连续性方法，即让 $f$ 从零形变到目标函数，同时让度量在 Kähler 度量空间中形变。把 (5.9) 看作从势到密度函数的映射，我们需要证明该映射是\textbf{淹没}且\textbf{真}的。淹没性由 Laplacian 上的标准 Hodge 理论推出。为证明真性，需要对 $\phi$ 的\textbf{先验估计}，且估计只依赖于 $X$ 与 $f$ 的界；因此我们可以对 $f$ 施加任意阶的正则性。

第一步是估计 $\|\phi\|_{C^0}$（称为“$C^0$-估计”），所用技巧称为 \textbf{Moser 迭代}。我们把方程改写成
\[
(1-e^f)\omega^n = -dd^c\phi\wedge(\omega^{n-1}+\omega^{n-1}\wedge\omega_\phi+\dots+\omega_\phi^{n-1}).
\]

% [Tutor Note: 这个恒等式来自 $\omega^n-\omega_\phi^n=-\sum_{k=0}^{n-1}dd^c\phi\wedge\omega^{n-1-k}\wedge\omega_\phi^{k}$ 的展开：括号中是 $n$ 项，首项 $\omega^{n-1}$、末项 $\omega_\phi^{n-1}$，中间项的 $\omega$ 与 $\omega_\phi$ 的幂次互补，容易多写或漏写一项。]
